%! Author = Hasan Nahiyan Nobel

% Preamble
\documentclass[11pt, a4paper]{article}
\title{There's a Blue Light}
\author{Hasan Nahiyan Nobel}
\date{October 2024}

% Packages
\usepackage{graphicx}
\graphicspath{{images/}}

% Document
\begin{document}
    % Show the title
    \maketitle
    % Show the TOC
    \tableofcontents


    \section{Introduction}\label{sec:introduction}
    This is just some normal text.


    \section{Images}\label{sec:images}
    And here is an image:
    \begin{figure}[h]
        \centering
        \includegraphics[width=0.75\textwidth]{latex_project_logo}
        \caption{The Latex Project Logo}
        \label{fig:latex_project_logo}
    \end{figure}


    \section{Lists}\label{sec:lists}
    This is an unordered list:
    \begin{itemize}
        \item The individual entries are indicated with a black dot, a so-called bullet.
        \item The text in the entries may be of any length.
    \end{itemize}
    And this list is ordered:
    \begin{enumerate}
        \item This is the first entry in our list.
        \item The list numbers increase with each entry we add.
    \end{enumerate}


    \section{Math}\label{sec:math}
    About math, say, $E = mc^2$ is an inline math.
    And the following is called ‘display math’ mode:
    \begin{equation}
        E=mc^2
        \label{eq:equation_1}
    \end{equation}


    \section{Table}\label{sec:table}
    This should be a basic table:
    \begin{table}[h]
        \centering
        \begin{tabular}{| c | c | c | c|}
            \hline
            Col1 & Col2 & Col2  & Col3 \\
            \hline
            1    & 6    & 87837 & 787  \\
            \hline
            2    & 7    & 78    & 5415 \\
            \hline
            3    & 545  & 778   & 7507 \\
            \hline
            4    & 545  & 18744 & 7560 \\
            \hline
            5    & 88   & 788   & 6344 \\
            \hline
        \end{tabular}
        \label{tab:table}
    \end{table}

\end{document}